% !TEX program = xelatex
\documentclass[11pt,a4paper]{article}
\usepackage{fontspec}
\setmainfont{Calibri}
\setmonofont[Scale=0.88]{Consolas}
\newfontfamily\symbolfont{Segoe UI Symbol}
\usepackage{polyglossia}\setdefaultlanguage{polish}
\usepackage[table]{xcolor}
\usepackage{booktabs,longtable,array,geometry,enumitem,graphicx,fancyhdr,caption}
\usepackage[most]{tcolorbox}
\PassOptionsToPackage{hyphens}{url}
\usepackage[hidelinks,unicode]{hyperref}
\emergencystretch=2.5em
\geometry{a4paper,margin=2.1cm,top=2.3cm,bottom=2.3cm}
\setlength{\parindent}{0pt}\setlength{\parskip}{5pt}
\definecolor{apteal}{HTML}{0E7C86}
\definecolor{apteallight}{HTML}{E6F4F5}
\definecolor{apgrey}{HTML}{F3F4F6}
\definecolor{apamber}{HTML}{B45309}
\definecolor{apamberlight}{HTML}{FEF3C7}
\definecolor{apred}{HTML}{B91C1C}
\definecolor{apredlight}{HTML}{FEE2E2}
\definecolor{apgreen}{HTML}{15803D}
\definecolor{apgreenlight}{HTML}{DCFCE7}
\renewcommand{\arraystretch}{1.25}
\captionsetup{font=small,labelfont=bf}
\pagestyle{fancy}\fancyhf{}
\fancyfoot[L]{\small AMPERE POINT — testy bez Wi-Fi · Q11 z modułem Shelly · v1 · 2026-09-28}
\fancyfoot[R]{\small\thepage}
\renewcommand{\headrulewidth}{0pt}
\newtcolorbox{uwaga}[1][Uwaga]{enhanced,breakable,colback=apamberlight,colframe=apamber,title={#1},fonttitle=\bfseries,boxrule=0.6pt,arc=2pt,left=6pt,right=6pt,top=4pt,bottom=4pt}
\newtcolorbox{info}[1][Jak to działa]{enhanced,breakable,colback=apteallight,colframe=apteal,title={#1},fonttitle=\bfseries,boxrule=0.6pt,arc=2pt,left=6pt,right=6pt,top=4pt,bottom=4pt}
\newtcolorbox{wazne}[1][Ważne]{enhanced,breakable,colback=apredlight,colframe=apred,title={#1},fonttitle=\bfseries,boxrule=0.6pt,arc=2pt,left=6pt,right=6pt,top=4pt,bottom=4pt}
\newtcolorbox{przyklad}[1][Przykład liczbowy]{enhanced,breakable,colback=apgreenlight,colframe=apgreen,title={#1},fonttitle=\bfseries,boxrule=0.6pt,arc=2pt,left=6pt,right=6pt,top=4pt,bottom=4pt}
\newtcolorbox{kodbox}{enhanced,breakable,colback=apgrey,colframe=apgrey,boxrule=0pt,arc=2pt,left=5pt,right=5pt,top=3pt,bottom=3pt,fontupper=\ttfamily\footnotesize}
\setlist{nosep,leftmargin=*}
\setcounter{secnumdepth}{2}
\begin{document}

{\LARGE\bfseries Testy bez Wi-Fi: Q11 z modułem Shelly\par}
\vspace{4pt}\textcolor{apteal}{\rule{\textwidth}{1.2pt}}
\emph{Poziom D: nie ma routera; zostaje Bluetooth albo sieć modułu i telefon obok ładowarki. Do tego funkcje samego sterownika. Testy D1–D9 · AMPERE POINT · 28 września 2026 · oprac. Dawid Cekała}


\section*{W skrócie}

\begin{itemize}
\item \textbf{Sytuacja z życia:} garaż bez zasięgu Wi-Fi. Q11 na Tuya ma wtedy tylko menu ładowarki.
\item \textbf{Q11 na Shelly ma dwie drogi:}
\begin{itemize}
\item Bluetooth: aplikacja albo laptop;
\item własna sieć modułu, do której telefon łączy się bezpośrednio i otwiera stronę WWW modułu.
\end{itemize}
\item \textbf{Kłopot to godzina.} Bez Wi-Fi nie ma serwera czasu, a czasu ustawionego ręcznie moduł sterownikowi nie przekazuje (sprawdzone 25.09). D5 sprawdza, czy strona albo aplikacja w ogóle ustawiają zegar modułu. D6 sprawdza, co sterownik robi z oknem godzin bez godziny.
\item \textbf{Funkcje samego sterownika z menu} (D7, D9) też są w tym dokumencie, bo bez Wi-Fi to jedyne, co zostaje po stronie Tuya.
\item \textbf{Większość kroków robimy z modułem nadal w sieci biura,} żeby widzieć log. Prawdziwe wyłączenie Wi-Fi modułu (D8) robimy dopiero, gdy Bluetooth i strona WWW działają.
\item Poziomy łączności, pojęcia i porównanie z Tuya są w dokumencie 00, w tabeli 3.3.
\end{itemize}

\section{Przygotowanie}

\begin{itemize}
\item Laptop: karta Bluetooth (Intel), Python z biblioteką \texttt{bleak}.
\item \textbf{Narzędzie do napisania:} \texttt{narzedzia\textbackslash{}ble\_\allowbreak{}rpc.\allowbreak{}py}, klient poleceń Shelly przez Bluetooth. Ta sama treść JSON co przez Wi-Fi, opakowana w ramkę Bluetooth Shelly.
\item Telefon z aplikacją Shelly.
\item Monitor logu przez Wi-Fi uruchomiony (poza D8).
\item Symulator auta podłączony (D6–D8).
\end{itemize}

\section{Testy}

\subsection{D1. Skan Bluetooth}

\begin{longtable}{>{\raggedright\arraybackslash}p{3.00cm}>{\raggedright\arraybackslash}p{4.18cm}>{\raggedright\arraybackslash}p{8.10cm}}
\toprule
\rowcolor{apteallight}\textbf{Krok} & \textbf{Co robimy} & \textbf{Wynik oczekiwany} \\
\midrule\endhead
\bottomrule\endfoot
D1a & Skan Bluetooth z laptopa & moduł widoczny pod swoją nazwą, z reklamą Shelly \\
\end{longtable}

\subsection{D2. Polecenia z laptopa przez Bluetooth}

\begin{longtable}{>{\raggedright\arraybackslash}p{3.00cm}>{\raggedright\arraybackslash}p{5.40cm}>{\raggedright\arraybackslash}p{6.88cm}}
\toprule
\rowcolor{apteallight}\textbf{Krok} & \textbf{Co robimy} & \textbf{Wynik oczekiwany} \\
\midrule\endhead
\bottomrule\endfoot
D2a & Odczyt stanu przez Bluetooth & te same pola co przez Wi-Fi \\
D2b & Zmiana limitu prądu na 10 A przez Bluetooth & sterownik potwierdza w logu w ciągu 1 s, tak jak przy Wi-Fi \\
\end{longtable}

\subsection{D3. Aplikacja Shelly przez Bluetooth}

\begin{longtable}{>{\raggedright\arraybackslash}p{3.00cm}>{\raggedright\arraybackslash}p{5.96cm}>{\raggedright\arraybackslash}p{6.32cm}}
\toprule
\rowcolor{apteallight}\textbf{Krok} & \textbf{Co robimy} & \textbf{Wynik oczekiwany} \\
\midrule\endhead
\bottomrule\endfoot
D3a & Telefon z wyłączonym Wi-Fi i danymi komórkowymi; aplikacja Shelly: włącz/wyłącz ładowanie, zmiana limitu & \textbf{do ustalenia.} Aplikacja wymaga konta w chmurze, a chmura modułu jest wyłączona; może nie pokazać urządzenia. Zapisujemy, co widać i ile trwa połączenie \\
\end{longtable}

\subsection{D4. Strona WWW przez sieć modułu}

\begin{longtable}{>{\raggedright\arraybackslash}p{3.00cm}>{\raggedright\arraybackslash}p{8.35cm}>{\raggedright\arraybackslash}p{3.93cm}}
\toprule
\rowcolor{apteallight}\textbf{Krok} & \textbf{Co robimy} & \textbf{Wynik oczekiwany} \\
\midrule\endhead
\bottomrule\endfoot
D4a & Sprawdzamy ustawienia sieci modułu: czy jest włączona i \textbf{czy ma hasło} & zapisujemy \\
D4b & Telefon łączy się z siecią modułu; przeglądarka otwiera stronę modułu & strona działa bez routera i bez aplikacji \\
D4c & Zmiana limitu i włącznika ze strony & sterownik potwierdza \\
\end{longtable}

\textbf{Wniosek do analizy GAP, niezależnie od wyniku.} Jeśli sieć modułu nie ma hasła, każdy w pobliżu może sterować ładowarką przez stronę WWW.


\subsection{D5. Czas z telefonu}

\textbf{Cel.} Bez Wi-Fi jedynym źródłem godziny jest telefon. Sprawdzić, czy strona WWW albo aplikacja ustawiają zegar modułu przy połączeniu i czy sterownik dostaje tę godzinę. Wiemy już, że czasu ustawionego poleceniem moduł sterownikowi nie przekazuje (25.09).


\begin{longtable}{>{\raggedright\arraybackslash}p{3.00cm}>{\raggedright\arraybackslash}p{6.14cm}>{\raggedright\arraybackslash}p{6.14cm}}
\toprule
\rowcolor{apteallight}\textbf{Krok} & \textbf{Co robimy} & \textbf{Wynik oczekiwany / co zapisujemy} \\
\midrule\endhead
\bottomrule\endfoot
D5a & Zatrzymujemy kontener z serwerem czasu, restartujemy moduł & moduł nie ma godziny; sterownik dostaje zerową datę \\
D5b & Telefon: strona WWW modułu przez sieć modułu & czy zegar modułu dostał godzinę z telefonu; co moduł wysyła sterownikowi \\
D5c & Telefon: aplikacja przez Bluetooth & to samo \\
D5d & Serwer czasu z powrotem, restart modułu & poprawna godzina wraca \\
\end{longtable}

\textbf{Jeśli zegar modułu dostaje godzinę, a sterownik nie,} to potwierdza prośbę do Shelly: czas ustawiony z telefonu ma się liczyć jako ważny dla sterownika.


\subsection{D6. Okno godzin bez godziny}

\textbf{Cel.} Co sterownik robi z oknem godzin, gdy dostaje od modułu zerową datę. Tak będzie po każdym zaniku zasilania bez Wi-Fi.


\begin{longtable}{>{\raggedright\arraybackslash}p{3.00cm}>{\raggedright\arraybackslash}p{7.23cm}>{\raggedright\arraybackslash}p{5.05cm}}
\toprule
\rowcolor{apteallight}\textbf{Krok} & \textbf{Co robimy} & \textbf{Wynik oczekiwany / co zapisujemy} \\
\midrule\endhead
\bottomrule\endfoot
D6a & Serwer czasu zatrzymany, restart modułu; z menu tryb „harmonogram”, okno obejmujące bieżącą godzinę; symulator „ładuje” & czy sterownik ładuje, nie ładuje, czy robi coś innego \\
D6b & Okno poza bieżącą godziną & to samo \\
D6c & Serwer czasu z powrotem, restart modułu & sterownik wraca do normalnej pracy z oknem \\
\end{longtable}

\subsection{D7. Funkcje sterownika z menu}

\textbf{Cel.} Te funkcje działają bez żadnej łączności. Sprawdzamy, czy sterownik zgłasza zmiany z menu modułowi, żeby moduł je pokazał.


\begin{longtable}{>{\raggedright\arraybackslash}p{3.00cm}>{\raggedright\arraybackslash}p{4.70cm}>{\raggedright\arraybackslash}p{7.58cm}}
\toprule
\rowcolor{apteallight}\textbf{Krok} & \textbf{Co robimy} & \textbf{Wynik oczekiwany} \\
\midrule\endhead
\bottomrule\endfoot
D7a & Limit prądu z menu ładowarki: 8 A & sterownik sam zgłasza 8 A; pole w module aktualizuje się bez naszego zapisu \\
D7b & Z menu: czy okno godzin da się ustawić z minutami & jeśli tak, założenie „tylko pełne godziny” jest błędne; poprawić w zestawieniu punktów danych \\
D7c & Start ładowania po podłączeniu auta, bez żadnego polecenia & jak 25.09: sterownik startuje sam; potwierdzenie na obecnej konfiguracji \\
\end{longtable}

\subsection{D8. Moduł naprawdę bez Wi-Fi}

\textbf{Cel.} Wszystko z D2–D4 przy wyłączonej stacji Wi-Fi modułu, czyli naprawdę bez routera.


\begin{longtable}{>{\raggedright\arraybackslash}p{3.00cm}>{\raggedright\arraybackslash}p{4.84cm}>{\raggedright\arraybackslash}p{7.44cm}}
\toprule
\rowcolor{apteallight}\textbf{Krok} & \textbf{Co robimy} & \textbf{Wynik oczekiwany} \\
\midrule\endhead
\bottomrule\endfoot
D8a & Sieć modułu włączona (D4a); harmonogram w module: za 3 min limit 6 A & droga awaryjna gotowa \\
D8b & Przez Bluetooth wyłączamy stację Wi-Fi modułu & moduł znika z sieci biura; log niedostępny \\
D8c & Powtórka D2, D3 i D4 & działa bez Wi-Fi; dowodem są wyświetlacz ładowarki i odpowiedzi przez Bluetooth \\
D8d & Harmonogram z D8a & czy odpalił; wyświetlacz pokazuje 6 A \\
D8e & Ikona sieci na wyświetlaczu & czy moduł wysyła inny stan niż „4”, np. „2” = brak routera \\
D8f & Przez Bluetooth włączamy stację Wi-Fi z powrotem & moduł wraca do sieci biura w ciągu minuty, log wraca. Jeśli moduł zażąda hasła sieci, wpisuje je użytkownik \\
\end{longtable}

\textbf{Ryzyko: średnie.} Jeśli Bluetooth zawiedzie, moduł zostaje bez Wi-Fi. Droga powrotu: sieć modułu z telefonu. Dlatego D8 dopiero po udanych D2 i D4.


\subsection{D9. Reset sieci z menu ładowarki}

\textbf{Cel.} W Q11 na Tuya pozycja menu „reset sieci” wprowadza moduł w parowanie. Sprawdzić, co to polecenie sterownika robi z modułem Shelly.


\begin{longtable}{>{\raggedright\arraybackslash}p{3.00cm}>{\raggedright\arraybackslash}p{5.55cm}>{\raggedright\arraybackslash}p{6.73cm}}
\toprule
\rowcolor{apteallight}\textbf{Krok} & \textbf{Co robimy} & \textbf{Wynik oczekiwany} \\
\midrule\endhead
\bottomrule\endfoot
D9a & Monitor na najwyższym poziomie logu; z menu ładowarki: reset sieci & w logu ramka od sterownika z poleceniem resetu Wi-Fi i reakcja modułu: ignoruje, kasuje Wi-Fi albo wchodzi w parowanie \\
D9b & Jeśli moduł skasował Wi-Fi: ponowne dodanie do sieci & moduł wraca do sieci \\
\end{longtable}

\textbf{Drogi powrotu w D9b, w tej kolejności:}

\begin{enumerate}
\item \textbf{Aplikacja Shelly przez Bluetooth.} Uwaga: 24.09 ta droga zawiodła, bo aplikacja zapisała nazwę sieci ze spacją na końcu.
\item \textbf{Sieć modułu i jego strona WWW.} Nazwę sieci i hasło wpisuje użytkownik.
\end{enumerate}

USB odpada, bo moduł jest zasilany z płytki.


\textbf{Ryzyko: duże.} Utrata Wi-Fi modułu, dlatego to ostatni test całego programu. Wynik zapisujemy jako fakt do dokumentacji dla fabryki.


\section{Kolejność i czas}

\begin{longtable}{>{\raggedright\arraybackslash}p{3.13cm}>{\raggedright\arraybackslash}p{5.99cm}>{\raggedright\arraybackslash}p{6.16cm}}
\toprule
\rowcolor{apteallight}\textbf{Kolejność} & \textbf{Testy} & \textbf{Czas} \\
\midrule\endhead
\bottomrule\endfoot
1 & narzędzie Bluetooth, D1–D4 & 1 h 35 min \\
2 & D5 & 15 min \\
3 & D6, D7 & 50 min \\
4 & D8 & 30 min \\
5 & D9, na końcu całego programu & 20 min plus ewentualne dodanie do sieci \\
\end{longtable}

Razem około 3 h 30 min. Miejsce tych testów w kolejności całego programu: dokument 00, rozdział 5.


\section{Wyniki}

\begin{longtable}{>{\raggedright\arraybackslash}p{7.73cm}>{\raggedright\arraybackslash}p{2.29cm}>{\raggedright\arraybackslash}p{2.52cm}>{\raggedright\arraybackslash}p{2.29cm}}
\toprule
\rowcolor{apteallight}\textbf{Test} & \textbf{Wynik} & \textbf{Co zaobserwowano} & \textbf{Data} \\
\midrule\endhead
\bottomrule\endfoot
D1 skan Bluetooth &  &  &  \\
D2 polecenia z laptopa &  &  &  \\
D3 aplikacja przez Bluetooth &  &  &  \\
D4 strona WWW przez sieć modułu &  &  &  \\
D5 czas z telefonu &  &  &  \\
D6 okno godzin bez godziny &  &  &  \\
D7 funkcje z menu &  &  &  \\
D8 moduł bez Wi-Fi &  &  &  \\
D9 reset sieci z menu &  &  &  \\
\end{longtable}

\end{document}